% ===================================================================
%  সারণি নং 10 — সমুদ্র। — বন্দর।
%
%  Ceci n'est PAS une transcription : c'est une traduction, faite en
%  2026, en bengali d'aujourd'hui. Voir texto/bn/10-sarani-01.tex
%  pour la consigne, qui vaut pour les seize fichiers.
% ===================================================================

%%K t10-apar-1 apar
\VUsaut{4.0mm}
\VUtitre{15.0pt}{60}{সারণি নং 10}
\VUsaut{3.0mm}
\VUpk{11.4pt}{40}{সমুদ্র। --- বন্দর।}
\VUsaut{3.0mm}

%%K t10-01-1 p
1. --- গরমকালে যখন কেউ কেউ পাহাড়ে শান্তি ও ঠান্ডা হাওয়া খোঁজে, অন্যেরা
সমুদ্রতীরে যেতে বেশি পছন্দ করে। গত বছর আমরা তা-ই করেছিলাম, কারণ আমাদের
\VUgras{মহাসাগর} \textsuperscript{(1)} খুব ভালো লাগে।

%%K t10-02-1 p
2. --- আমার তো \VUgras{দিগন্ত} \textsuperscript{(2)} পর্যন্ত ছড়ানো সেই বিশাল
জলচাদর দেখতে বড় সুখ হয়, আর এটি দেখতে যে বিশাল \VUgras{বাষ্পজাহাজ}
\textsuperscript{(3)} আমেরিকাগামী যাত্রীদের নিয়ে দীর্ঘ \VUgras{সমুদ্রযাত্রায়}
বেরোয়।

%%K t10-03-1 p
3. --- তাই আমার বাবা, যিনি জানেন আমার কী ভালো লাগে, ছুটির জন্য সব সময় আমাদের
কোনো বড় \VUgras{নৌ-বন্দরের} কাছের খুব শান্ত একটি তীর বাছেন।

%%K t10-03-2 p
আমাদের গত প্রবাসে \VUgras{নৌবহর} এসে সেই বিশাল \VUgras{ঢেউ-রোধকের}
\textsuperscript{(4)} পিছনে নোঙর ফেলেছিল যা \VUgras{নৌ-বন্দরকে}
\textsuperscript{(5)} বন্ধ করে ও তার রক্ষা করে, আর আমি নিশ্চিন্তে আলাদা আলাদা
\VUgras{রণতরি} দেখতে পেলাম --- ছোট \VUgras{ডুবোজাহাজ} \textsuperscript{(10)}
থেকে, যা ডুব দিয়ে জলের নিচে মিলিয়ে যায়, বিশাল \VUgras{রণপোত}
\textsuperscript{(9)} পর্যন্ত, যাকে এ পর্যন্ত \VUgras{টর্পেডো নৌকার}
\textsuperscript{(8)} হামলা ছাড়া কিছু ভয় দেখাত না।

%%K t10-tit-2 sub
\VUsaut{3.0mm}
\VUpk{10.2pt}{40}{ছোট নৌকা।}
\VUsaut{3.0mm}

%%K t10-04-1 p
4. --- সে আগেই \VUgras{চতুরতায়} মোটা ধাতব বর্মের দুর্বল জায়গাটি খুঁজে নিত ও
সেখানে সেই বিপজ্জনক \VUgras{টর্পেডো} গেঁথে দিত যার বিস্ফোরণে জাহাজ ডুবে যায়;
কিন্তু সে ডুবোজাহাজের চেয়ে তবু কম ভয়ংকর ছিল, কারণ ধ্বংসকারীরা সহজেই তার পিছু
নিতে পারে।

%%K t10-05-1 p
5. --- আমি দ্রুতগামী বর্মআঁটা \VUgras{ক্রুজারও} \textsuperscript{(6)} দেখতে
পেলাম, যারা সত্যিকারের রণপোতের মতোই অভেদ্য, কয়েকটি \VUgras{সৈন্যবাহী জাহাজ}
\textsuperscript{(12)}, তিন-চারটি \VUgras{কামান-নৌকা} \textsuperscript{(7)}, আর
শেষে একটি \VUgras{ফ্রিগেট} \textsuperscript{(11)}। \VUgras{সেই নৌবহরকে নোঙর
তুলতে দেখাই সবচেয়ে কৌতূহলজনক।}

%%K t10-06-1 p
6. --- ইস্পাতের সেই বিশাল স্তূপ আর সেই সুরুচিসম্মত \VUgras{ত্রিমাস্তুল} বা
সুঠাম \VUgras{পালতোলা জাহাজের} \textsuperscript{(13)} গড়নে কত তফাত, যা এখনও
বাণিজ্যবহরে কাজে লাগে এবং যাদের আমি \VUgras{ঘাটবাঁধে} \textsuperscript{(14)}
বেড়াতে বেড়াতে পুরো পাল খুলে বন্দরে ঢুকতে দেখতাম। এ কথা সত্য যে
\VUgras{হাওয়া} বিপরীত হলে তারা নিজেদের অনেক সুবিধা হারাত। ডকে ঢোকার জন্য
তাদের সব পাল গুটাতে হত, আর একটি \VUgras{টানা-নৌকা} \textsuperscript{(16)} লাগত
যা তাদের আনতে যাবে ও সাধারণ \VUgras{মালবাহী নৌকার} \textsuperscript{(17)} মতো
ঘাট পর্যন্ত নিয়ে আসবে।

%%K t10-tit-3 sub
\VUsaut{3.0mm}
\VUpk{10.2pt}{40}{বাণিজ্যবন্দর।}
\VUsaut{3.0mm}

%%K t10-07-1 p
7. --- যে বন্দরের কথা আমি বলছি তা এইজন্য এত কৌতূহলজনক যে তাতে কেবল বিশাল
কাজে কৃত্রিমভাবে গড়া নৌ-বন্দরই নয়, বরং একটি বড় নদীর \VUgras{মোহনায়} একটি
অপূর্ব \VUgras{বাণিজ্যবন্দরও} আছে।

%%K t10-08-1 p
8. --- দুই ডককে আলাদা করা ভূমির ফালিটি দুর্গপ্রাচীরে ঘেরা এবং সমুদ্র পর্যন্ত
ছড়ানো \VUgras{কামানের সারি} ও \VUgras{বুরুজে} রক্ষিত। \VUgras{প্রাচীরে}
\textsuperscript{(15)} বেড়ানোর জন্য বিশেষ অনুমতি লাগে। আমার তা নিজের কাকার
মাধ্যমে সহজেই মিলে গিয়েছিল, যিনি \VUgras{জাহাজ-নির্মাণশালার}
\textsuperscript{(18)} প্রকৌশলী, আর আমি বিশাল ছাউনির নিচে গড়ে ওঠা জাহাজ দেখতে
গেলাম।

%%K t10-09-1 p
9. --- আমি একটি রণপোতকে জলে নামতেও দেখেছি, আর আমার সেই রোমাঞ্চ মনে আছে যা
সারা ভিড় অনুভব করেছিল যখন সেই বিশাল স্তূপ, নিজের ভরসায় ছেড়ে দেওয়া, নিজের
পাটাতনে পিছলে নামতে লাগল ও নদীর জলে ভাসতে এল।

%%K t10-10-1 p
10. --- নির্মাণশালায় পৌঁছাতে \VUgras{কুচকাওয়াজের মাঠ} \textsuperscript{(19)}
পেরোতে হয়, যেখানে ছাউনির সেনা প্রতিদিন কসরত করতে আসে, আর লোহার সেই
\VUgras{সাঁকো} \textsuperscript{(20)} পেরোতে হয় যা কুচকাওয়াজের মাঠকে
\VUgras{অস্ত্রাগারের} \textsuperscript{(21)} সঙ্গে জোড়ে।

%%K t10-tit-4 sub
\VUsaut{3.0mm}
\VUpk{10.2pt}{40}{অস্ত্রাগার ও ডক।}
\VUsaut{3.0mm}

%%K t10-11-1 p
11. --- নিজের কাকার সঙ্গে আমি সেই বড় প্রতিষ্ঠান দেখলাম, যা \VUgras{কামানে}
\textsuperscript{(22)}, \VUgras{গোলায়} \textsuperscript{(23)} ও \VUgras{খোলে}
ভরা। আমি সেই \VUgras{লোহশালাও} \textsuperscript{(24)} দেখলাম যেখানে বাষ্পজাহাজের
\VUgras{বয়লার} \textsuperscript{(25)} তৈরি হয়।

%%K t10-12-1 p
12. --- খুব কাছেই \VUgras{ডক} \textsuperscript{(26)} --- শুকনো ডক --- আর সেই খুব
উল্লেখযোগ্য \VUgras{ভাসমান ডক} \textsuperscript{(27)}, যেখানে আমি মেরামত হতে থাকা
একটি জাহাজে কাছ থেকে কোনো পোতের \VUgras{পাখা} \textsuperscript{(28)},
\VUgras{তলি} \textsuperscript{(29)} ও \VUgras{হাল} \textsuperscript{(30)} দেখতে
পেলাম।

%%K t10-13-1 p
13. --- বাণিজ্যবন্দর নৌ-বন্দরের মতোই দেখার যোগ্য, আর আমি অনেক ঘণ্টা সেই ঘাটে
কাটিয়েছি যেখানে নদী বেয়ে ওঠা \VUgras{চাকাওয়ালা বাষ্পজাহাজ}
\textsuperscript{(31)} বাঁধা থাকে, আর সেই \VUgras{তেপায়া ক্রেনটিকে}
\textsuperscript{(32)} কাজ করতে দেখেছি, যা বিশাল বোঝা পালকের মতো তুলে নেয়।

%%K t10-tit-5 sub
\VUsaut{4.0mm}
\VUpk{10.2pt}{40}{কর্ণধার।}
\VUsaut{3.0mm}

%%K t10-14-1 p
14. --- আমি সেই \VUgras{আন্তঃমহাসাগরীয় জাহাজগুলি} \textsuperscript{(34)}
নিহারলাম যা নিজেদের \VUgras{পতাকা} \textsuperscript{(35)} উড়িয়ে বন্দর ছাড়ে,
যাদের একজন দক্ষ \VUgras{কর্ণধার} \textsuperscript{(36)} চালান, যিনি
\VUgras{বয়ায়} \textsuperscript{(40)} ও \VUgras{সংকেতস্তম্ভে} চিহ্নিত
\VUgras{জলপথে} অপূর্বভাবে চলতে জানেন, আর সেই \VUgras{মালনৌকা}
\textsuperscript{(41)} এড়িয়ে যান যা নিজেদের একমাত্র চৌকো পালে বড় কষ্টে চালানো
হয়। আমি \VUgras{গলুইয়ে} \textsuperscript{(37)} দ্বিতীয় শ্রেণির
\VUgras{কেবিন} \textsuperscript{(39)} চিনতে পারতাম, আর \VUgras{পিছনের ভাগে}
\textsuperscript{(38)} প্রথম শ্রেণির যাত্রীদের কক্ষ।

%%K t10-15-1 p
15. --- কখনও কখনও আমি কোনো \VUgras{মালবাহী জাহাজের} \textsuperscript{(42)}
বোঝাইও দেখতাম, যাতে \VUgras{গুদাম} \textsuperscript{(43)} ও \VUgras{বন্দর-স্টেশনের}
\textsuperscript{(44)} মাল ঠাসা হত, যা \VUgras{ঘাটের রেললাইনের}
\textsuperscript{(45)} অনেক গাড়ি আনত।

%%K t10-tit-6 sub
\VUsaut{3.0mm}
\VUpk{10.2pt}{40}{নৌকাবিহার।}
\VUsaut{3.0mm}

%%K t10-16-1 p
16. --- আমি প্রায়ই \VUgras{ভাসমান ডকে} \textsuperscript{(53)} নৌকা চালাতাম,
যেখানে ছোট \VUgras{নৌকারা} \textsuperscript{(46)} আশ্রয় নেয়, আর
\VUgras{বৈঠা} \textsuperscript{(47)} চালানো আমার কাছে আনন্দ ছিল। আমি
\VUgras{পাটাতনে} \textsuperscript{(48)} উঠে যেতাম, কোনো \VUgras{পালতোলা নৌকার}
\textsuperscript{(49)}, \VUgras{দড়াদড়ির} \textsuperscript{(52)} সাহায্যে
\VUgras{মাস্তুলে} \textsuperscript{(51)} \VUgras{আড়াআড়ি দণ্ড}
\textsuperscript{(50)} পর্যন্ত চড়তাম, তারপর একটি \VUgras{ডিঙিতে}
\textsuperscript{(54)} নিজের বেড়ানো চালিয়ে যেতাম, ভারী \VUgras{মাছধরা নৌকার}
\textsuperscript{(55)} মাঝে, যেখানে \VUgras{জাল} \textsuperscript{(56)} শুকোতে
ঝোলানো থাকে।

%%K t10-17-1 p
17. --- \VUgras{ঘাটে} \textsuperscript{(58)} আমার সামনে দিয়ে নানা রকম
\VUgras{মাল} \textsuperscript{(59)} যেত, \VUgras{বাক্সে} \textsuperscript{(60)} বা
\VUgras{বোঝায়} \textsuperscript{(61)} বাঁধা। সেগুলি মালনৌকার \VUgras{লদান}
\textsuperscript{(63)} গড়ত, আর শক্তিশালী \VUgras{ক্রেন} \textsuperscript{(62)}
সেগুলি তুলে \VUgras{ট্রাকে} \textsuperscript{(64)} রাখত, যতক্ষণ
\VUgras{বহনকারীরা} \VUgras{শুল্কঘরে} \textsuperscript{(65)} সেগুলির ঘোষণা করত ও
শুল্ক মেটাত।

%%K t10-18-1 p
18. --- শেষে, \VUgras{ঘোরানো সেতু} \textsuperscript{(66)} বা \VUgras{জলকপাট}
\textsuperscript{(69)} পর্যন্ত পৌঁছে, সেই \VUgras{খননকারী নৌকার}
\textsuperscript{(68)} পাশ দিয়ে যেতে যেতে যা \VUgras{খাল} \textsuperscript{(67)}
সাফ করে, আমি \VUgras{সংকেত-মাস্তুলের} \textsuperscript{(70)} পাশে ঘাটে নেমে
পড়তাম।

%%K t10-19-1 p
19. --- সেই ঘাটের অন্য দিকে সেই \VUgras{ছোট নৌকাগুলি} \textsuperscript{(71)}
ভেড়ে যা নৌবহরের অফিসারদের কূলে আনে। নাবিকদের তালে তালে বৈঠা তুলতে দেখা সুন্দর
দৃশ্য, যখন \VUgras{ঢেউয়ে} \textsuperscript{(72)} ওঠা নৌকা সহজেই নামার সিঁড়ি
পর্যন্ত এসে লাগে। এখান থেকেই \VUgras{জোয়ার} আর \VUgras{তীরে}
\textsuperscript{(73)} \VUgras{ভাটা} ও \VUgras{জোয়ারের} চলন সবচেয়ে ভালো দেখা
যায়।

%%K t10-tit-7 sub
\VUsaut{3.0mm}
\VUpk{10.2pt}{40}{অফিসারদের ক্লাব।}
\VUsaut{3.0mm}

%%K t10-20-1 p
20. --- বাড়ি ফেরার জন্য আমি \VUgras{বৈদ্যুতিক ট্রাম} \textsuperscript{(74)}
নিতাম, যার \VUgras{থামার জায়গা} \textsuperscript{(75)} অফিসারদের ক্লাবের সামনে
দাঁড়ানো \VUgras{খুঁটিতে}।

%%K t10-20-2 p
ক্লাবের \VUgras{বারান্দা} \textsuperscript{(76)} থেকে, যা \VUgras{নোঙরে}
\textsuperscript{(77)}, কামানে ও গোলায় সাজানো, আমি প্রায়ই নৌ-অফিসারদের
চমৎকার জমায়েত দেখেছি: নিজেদের তলোয়ারসহ \VUgras{লেফটেন্যান্ট}
\textsuperscript{(78)}, \VUgras{গোলন্দাজ বাহিনীর} \textsuperscript{(79)} ও
\VUgras{পদাতিক বাহিনীর} \textsuperscript{(80)} \VUgras{ক্যাপ্টেন} নিজেদের
\VUgras{কৃপাণসহ}। তাঁদের পিছনে একজন \VUgras{নাবিক} \textsuperscript{(81)}
দাঁড়িয়ে থাকত, যে তাঁদের একটি \VUgras{দূরবিন} এনে দিত যখন তাঁরা দূরের জিনিস
চিনতে চাইতেন।

%%K t10-21-1 p
21. --- আমি তরুণ \VUgras{উপ-লেফটেন্যান্টদেরও} \textsuperscript{(82)} দেখেছি, যাঁরা
নিজেদের \VUgras{কমান্ডার} \textsuperscript{(83)} বা \VUgras{অ্যাডমিরালের}
\textsuperscript{(84)} কাছে আদেশ নিতেন, যতক্ষণ একজন \VUgras{কোয়ার্টারমাস্টার}
\textsuperscript{(85)} সেগুলি জাহাজে নিয়ে যেতে অপেক্ষা করত।

%%K t10-22-1 p
22. --- সেই বারান্দা থেকে দৃশ্য চমৎকার: সেখান থেকে সারা তীর দেখা যায়, নিজের
\VUgras{তাঁবু} \textsuperscript{(86)} ও \VUgras{স্নানের কুঁড়েসহ}
\textsuperscript{(87)}, আর স্নানের সময় \VUgras{স্নানকারীদের}
\textsuperscript{(88)} \VUgras{জুয়াঘরের} \textsuperscript{(89)} সামনে প্রসন্ন
লাফাতে দেখা যায়।

%%K t10-23-1 p
23. --- তীরের ডগায় কূলটি একটি ছোট পাথুরে \VUgras{উপদ্বীপ}
\textsuperscript{(91)} গড়ে, যেখানে \VUgras{ঢেউ} \textsuperscript{(92)} আছড়ে
ভাঙে এবং যার উপরে একটি \VUgras{বাতিঘর} \textsuperscript{(93)} গড়া, যা রাতে
নাবিকদের পথ দেখায়। এই উপদ্বীপ মূল ভূমির সঙ্গে কেবল একটি খুব সরু
\VUgras{যোজকে} \textsuperscript{(90)} জোড়া।

%%K t10-tit-8 sub
\VUsaut{3.0mm}
\VUpk{10.2pt}{40}{তীরের সাধারণ দৃশ্য।}
\VUsaut{3.0mm}

%%K t10-24-1 p
24. --- \VUgras{খাড়া পাড়} \textsuperscript{(94)} এগিয়ে চলে, সমুদ্রে কাটা, যা
তাতে \VUgras{উপসাগর} \textsuperscript{(95)} ও \VUgras{অন্তরীপের} সারি খোদাই
করে, আর তাকে খুব মনোরম করে তোলে।

%%K t10-24-2 p
\VUgras{মালভূমিতে} \textsuperscript{(96)}, যা \VUgras{প্রণালীর}
\textsuperscript{(98)} উপরে ঝুঁকে, একটি খুব গুরুত্বপূর্ণ \VUgras{গড়}
\textsuperscript{(97)} আছে ও কয়েকটি কুঠি।

%%K t10-25-1 p
25. --- সেই প্রণালী মূল ভূমি থেকে একটি খুব বিপজ্জনক \VUgras{দ্বীপকে}
\textsuperscript{(99)} আলাদা করে, যা \VUgras{পাথরে} \textsuperscript{(100)} ও
\VUgras{প্রবালপ্রাচীরে} ঘেরা, যেখানে নৌকা ও বড় জাহাজ পর্যন্ত কখনও কখনও ভেঙে
যায়। উদ্ধারের সমস্ত তৎপরতা ও \VUgras{উদ্ধার-নৌকার} দ্রুত পৌঁছানো সত্ত্বেও এই
\VUgras{জাহাজভঙ্গ} ভয়ংকর হয়, আর বিষুবের \VUgras{ঝড়ে} কয়েকটি জাহাজ নিজেদের
সব লোক ও মালসমেত ডুবে গেছে।

%%K t10-25-2 p
এই প্রতিবেশ সত্ত্বেও \VUgras{বন্দর} নাবিকদের খুব আনাগোনা দেখে।
\VUsaut{6.0mm}
\VUfilet{22mm}
